\chapter{Analiza Seiberga--Wittena: od lokalnego przekroju do niezmiennika}
\label{ch:tg-sw-analiza-uzup}

W poprzednim rozdziale opisaliśmy mechanizm konstrukcji przestrzeni moduli.
Pozostały jednak dwa rodzaje przejść, których nie należy utożsamiać:
rachunki liniowe na przestrzeniach spinorów oraz twierdzenia analityczne
zapewniające regularność, zwartość i istnienie ilorazu. W tym rozdziale
rozpiszemy przejścia między tymi poziomami. Wykorzystamy twierdzenia
o oszacowaniach eliptycznych, twierdzenie Sarda--Smale'a i zasadę
jednoznacznego przedłużenia, podając wyraźnie moment ich użycia.
Nie będziemy dowodzić od podstaw teorii eliptycznych operatorów Fredholma.

Przez cały rozdział $X$ oznacza \emph{zamkniętą, spójną, gładką,
zorientowaną czterorozmaitość}, $g$ --- ustaloną metrykę riemannowską,
a $\mathfrak s$ --- ustaloną strukturę $\mathrm{Spin}^c$.
Piszemy $L=\det S^+$, gdzie $S^\pm\to X$ są jej wiązkami spinorowymi.
Formy o wartościach w $i\mathbb R$ traktujemy jako rzeczywiste
przestrzenie wektorowe. Oznaczenia $\mathcal C_k$, $\mathcal G_{k+1}$
oraz $\mathfrak F_\eta$ zachowują znaczenie z poprzedniego rozdziału.

\section{Normalizacja mnożenia Clifforda i oszacowanie spinora}
\label{sec:tg-sw12-clifford}

Iloczyn hermitowski na $S^+$ przyjmujemy liniowy w pierwszym
argumencie. Ustalamy mnożenie Clifforda na całej wiązce
$S=S^+\oplus S^-$ przez relacje
\[
\rho(v)\rho(w)+\rho(w)\rho(v)=-2\langle v,w\rangle
\operatorname{id}_S,\qquad \rho(v)^*=-\rho(v)
\]
dla rzeczywistych kowektorów. Rozszerzamy je zespolono-liniowo
i na formy stopnia dwa przyjmujemy
\[
\rho(v\wedge w)=\tfrac12
\bigl(\rho(v)\rho(w)-\rho(w)\rho(v)\bigr).
\]
Nie wprowadzamy dodatkowego przeskalowania tej mapy.

W dodatnio zorientowanej bazie ortonormalnej
$(e^1,e^2,e^3,e^4)$ oznaczmy $c_j=\rho(e^j)$
i $V=c_1c_2c_3c_4$. Przyjmujemy, że $S^+$
jest przestrzenią własną $V$ dla wartości $-1$
(tj. wartość $+1$ zespolonego operatora objętości $-V$).
Z relacji Clifforda na $S^+$ otrzymujemy
$c_3c_4=c_1c_2$, $c_2c_4=-c_1c_3$ oraz
$c_2c_3=c_1c_4$. Na przykład
$c_1c_2 V=-c_3c_4$ i $V|_{S^+}=-\operatorname{id}$,
co daje pierwszą równość; pozostałe wynikają
z tych samych relacji antykomutacji.
Wobec tego dla rzeczywistej formy samodualnej
\[
\omega=x(e^{12}+e^{34})+
y(e^{13}-e^{24})+z(e^{14}+e^{23})
\]
mamy $\rho(i\omega)=2i(xc_1c_2+yc_1c_3+zc_1c_4)|_{S^+}$,
gdzie $e^{jk}=e^j\wedge e^k$.
Trzy operatory $c_1c_j$ ($j=2,3,4$)
są antyhermitowskie, mają kwadrat $-\operatorname{id}$
i parami antykomutują. Każdy ma ślad zero,
ponieważ sprzężenie przez inny z nich zmienia jego znak.
Wynika stąd, że $\rho(i\omega)$ jest hermitowski,
bezśladowy i
\[
\rho(i\omega)^2=4(x^2+y^2+z^2)\operatorname{id}_{S^+}
=2|\omega|^2\operatorname{id}_{S^+}.
\]
Zatem $\rho:i\Lambda^2_+\to\operatorname{Herm}_0(S^+)$
jest injektywny, a obie przestrzenie mają rzeczywisty
wymiar trzy, więc jest izomorfizmem.
Ponadto $\|\rho(i\omega)\|_{\mathrm{op}}=\sqrt2|\omega|$.
Formy antysamodualne działają na $S^+$ zerowo,
co również wynika z powyższych trzech równości.
Ta normalizacja jest zgodna z definicją
\ref{def:tg-sw-eqn}.

\begin{definicja}[Kwadratowe odwzorowanie spinorowe]
\label{def:tg-sw12-q}
Dla $x\in X$ i $\psi\in S_x^+$ niech
$P_\psi(v):=\langle v,\psi\rangle\psi$ oraz
$P_\psi^0:=P_\psi-\tfrac12|\psi|^2\operatorname{id}_{S_x^+}$.
Definiujemy $q(\psi)\in i\Lambda^2_{+,x}$ równością
$\rho(q(\psi))=P_\psi^0$.
Dla operatora Diraca i połączenia wyznacznikowego
przyjętych w definicji \ref{def:tg-sw-eqn}
stosujemy formułę Weitzenböcka--Lichnerowicza
\[
D_A^*D_A=\nabla_A^*\nabla_A+\frac{s_g}{4}
 +\frac12\rho(F_A^+),
\]
gdzie $s_g$ jest krzywizną skalarną. Czynnik $1/2$
przy krzywiźnie pochodzi z działania środkowego
okręgu struktury $\mathrm{Spin}^c$ z wagą jeden
na spinorach, a z wagą dwa na wyznaczniku.
Równanie krzywizny
zapisujemy jako $F_A^+=q(\psi)+i\eta$, dla rzeczywistej
samodualnej formy $\eta$.
\end{definicja}

W innych podręcznikach $q$ może różnić się czynnikiem $2$
lub znakiem, wraz z odpowiadającą mu konwencją Clifforda.
Przy porównywaniu formuł z różnych źródeł należy porównywać
\emph{całe równanie}, a nie sam symbol $q$.

\begin{lemat}[Dwie tożsamości algebraiczne]
\label{lem:tg-sw12-q}
Dla każdego $x\in X$ i $\psi,\phi\in S_x^+$ zachodzą
\[
\langle\rho(q(\psi))\psi,\psi\rangle=\frac12|\psi|^4,
\qquad
\rho((\mathrm dq)_\psi\phi)
=(\phi\otimes\psi^*+\psi\otimes\phi^*)_0.
\]
W szczególności $q(\psi)=0$ wtedy i tylko wtedy, gdy $\psi=0$.
\end{lemat}

\begin{proof}
Macierz operatora $P_\psi$ w bazie ortonormalnej, której
pierwszy wektor jest równoległy do niezerowego $\psi$,
ma wartości własne $|\psi|^2$ i $0$. Po odjęciu połowy
śladu otrzymujemy wartości $|\psi|^2/2$ i $-|\psi|^2/2$.
Stąd $P_\psi^0\psi=(|\psi|^2/2)\psi$ i pierwsza
równość. W przypadku $\psi=0$ obie strony znikają.
Dla drugiej równości różniczkujemy
$P_{\psi+t\phi}$ względem $t\in\mathbb R$ w zerze:
\[
\frac{\mathrm d}{\mathrm dt}\bigg|_{t=0}P_{\psi+t\phi}
=\phi\otimes\psi^*+\psi\otimes\phi^*.
\]
Przejście do części bezśladowej komutuje z różniczkowaniem.
Ponieważ $\rho$ jest izomorfizmem, otrzymujemy
drugą równość. Z pierwszej równości wynika
$q(\psi)=0\Rightarrow|\psi|^4/2=0$,
a implikacja odwrotna wynika z definicji.
Ponadto z obliczonej wcześniej normy Clifforda
i wartości własnych $P_\psi^0$ dostajemy
$|q(\psi)|=|\psi|^2/(2\sqrt2)$.
W szczególności oszacowanie w następnym
twierdzeniu można stosować z $C_\rho=\sqrt2$.
\end{proof}

\begin{twierdzenie}[Jawne oszacowanie typu Weitzenböcka]
\label{thm:tg-sw12-spinor-bound}
Niech $\eta\in\Omega^2_+(X)$ będzie gładka i niech $(A,\psi)$
będzie gładkim rozwiązaniem równań Seiberga--Wittena.
Niech $C_\rho>0$ spełnia
$\|\rho(i\omega)\|_{\mathrm{op}}\le C_\rho|\omega|$
dla każdej rzeczywistej formy samodualnej $\omega$.
Wtedy
\[
\|\psi\|_{L^\infty}^2\le
\|s_g^-\|_{L^\infty}+2C_\rho\|\eta\|_{L^\infty},
\qquad s_g^-:=\max\{0,-s_g\}.
\]
Jeżeli $s_g>0$ na całym $X$ i $\eta=0$, każde rozwiązanie
jest redukowalne.
\end{twierdzenie}

\begin{proof}
Z równania $D_A\psi=0$, tożsamości Weitzenböcka i
lematu \ref{lem:tg-sw12-q} otrzymujemy punktowo
\[
0=\frac12\Delta|\psi|^2+|\nabla_A\psi|^2
 +\frac{s_g}{4}|\psi|^2+\frac14|\psi|^4
 +\frac12\langle\rho(i\eta)\psi,\psi\rangle,
\]
gdzie $\Delta=\mathrm d^*\mathrm d$ na funkcjach.
W punkcie maksimum $|\psi|^2$ mamy $\Delta|\psi|^2\ge0$.
Odrzucając nieujemny wyraz $|\nabla_A\psi|^2$ oraz
szacując wyraz perturbacyjny z dołu przez
$-C_\rho|\eta||\psi|^2/2$, dostajemy
\[
\frac14|\psi|^4\le
\left(\frac{s_g^-}{4}+\frac{C_\rho}{2}|\eta|\right)
|\psi|^2.
\]
Jeżeli maksimum jest dodatnie, możemy podzielić przez
$|\psi|^2$; dla maksimum równego zero nierówność jest
oczywista. Gdy $s_g>0$ i $\eta=0$, w punkcie dodatniego
maksimum prawa strona pierwszej tożsamości byłaby dodatnia,
co daje sprzeczność. Zatem $\psi\equiv0$.
\end{proof}

\begin{przyklad}[Krzywizna dodatnia a niezmiennik]
Dla $X_k=\mathbb{CP}^2\#k\overline{\mathbb{CP}}^{\,2}$,
$0\le k\le8$, i $c=3H-\sum_{j=1}^kE_j$ mamy
$e(X_k)=3+k$, $\sigma(X_k)=1-k$ i $c^2=9-k>0$.
Metryka o dodatniej krzywiźnie skalarnej oraz $\eta=0$
wykluczają rozwiązania nieredukowalne. Rozwiązanie redukowalne
wymagałoby zarazem $c^+=0$, a wtedy $c^2\le0$.
Wobec tego w komorze zawierającej tę parę danych przestrzeń
moduli jest pusta. Nie jest to dowód zanikania niezmiennika
we wszystkich komorach, ponieważ $b_2^+(X_k)=1$.
\end{przyklad}

\section{Cechowanie: lokalny przekrój i globalna część harmoniczna}
\label{sec:tg-sw12-slice}

Lokalne ustalenie cechowania oraz globalne usunięcie części
dokładnej połączenia są różnymi operacjami. Pierwsza służy
budowaniu gładkiej struktury ilorazu w pobliżu rozwiązania
nieredukowalnego, druga --- dowodowi zwartości rodzin rozwiązań.

\begin{uwaga}[Zakres twierdzenia o przekroju]
Twierdzenie dotyczy konfiguracji nieredukowalnej oraz
małego otoczenia tożsamości w grupie cechowania.
W globalnym ilorazie występują tzw. duże transformacje
cechowania, klasyfikowane przez $H^1(X;\mathbb Z)$,
i nie można globalizować otrzymanej funkcji $f(\xi)$
bez dodatkowych warunków topologicznych.
Przy rozwiązaniach redukowalnych operator
$\mathrm d_0^*\mathrm d_0$ ma jądro funkcji stałych,
co uniemożliwia bezpośrednie zastosowanie tej samej
wersji twierdzenia o funkcji uwikłanej.
\end{uwaga}

\begin{twierdzenie}[Lokalny przekrój w punkcie nieredukowalnym]
\label{thm:tg-sw12-slice}
Ustalmy $k\ge3$ i gładką konfigurację
$\mathfrak c=(A,\psi)\in\mathcal C_k^*$.
Niech $\mathrm d_0$ oznacza różniczkę działania grupy
cechowania, zapisaną w poprzednim rozdziale, a
$\mathrm d_0^*$ --- jej formalne sprzężenie w rzeczywistym
iloczynie $L^2$. Istnieją otoczenie $U$ konfiguracji $\mathfrak c$ oraz
otoczenie $\mathcal U$ identyczności w grupie cechowania
takie, że dla każdej $\mathfrak c'\in U$ istnieje
dokładnie jedna transformacja $u\in\mathcal U$,
dla której $u\cdot\mathfrak c'\in
\mathfrak c+\ker\mathrm d_0^*$.
Jedyność dotyczy zatem małych transformacji
w ustalonym lokalnym modelu ilorazu.
\end{twierdzenie}

\begin{proof}
Dla $\mathfrak c=(A,\psi)$ ustalonej i małych
$\xi=(a,\phi)$ rozważamy funkcję
\[
\mathcal F(f,\xi):=\mathrm d_0^*
\bigl(e^{if}\cdot(\mathfrak c+\xi)-\mathfrak c\bigr),
\]
gdzie $f\in W^{k+1,2}(X;\mathbb R)$ jest mała, a $\xi$
ma klasę $W^{k,2}$. Działanie cechowania jest gładkie
na odpowiednich przestrzeniach Sobolewa, ponieważ $k\ge3$
i zachodzi właściwe twierdzenie o mnożeniu funkcji Sobolewa.
Mamy $\mathcal F(0,0)=0$, a różniczka względem $f$
w zerze jest równa $\mathrm d_0^*\mathrm d_0$.
Z lematu \ref{lem:tg-sw-full-coulomb} wiadomo, że
$\mathrm d_0^*\mathrm d_0=4\mathrm d^*\mathrm d+|\psi|^2$
ma zerowe jądro. Jest eliptyczny i samosprzężony o indeksie
zero, więc jest izomorfizmem
$W^{k+1,2}\to W^{k-1,2}$.
Twierdzenie o funkcji uwikłanej daje jedyną małą funkcję
$f=f(\xi)$ spełniającą $\mathcal F(f,\xi)=0$.
Jest to lokalna jednoznaczność względem małych
transformacji cechowania. Aby uzyskać mapę na ilorazie,
trzeba dodatkowo wykluczyć identyfikacje pomiędzy
punktami dostatecznie małego przekroju przez
transformacje spoza otoczenia identyczności.

W tym celu przypuśćmy, że $\mathfrak c_n,\mathfrak c'_n$
zbiegają do $\mathfrak c$ w $\mathcal C_k$ i
$u_n\cdot\mathfrak c_n=\mathfrak c'_n$.
Z równania połączeń wynika
$\mathrm du_n=u_n(A_n-A'_n)/2$.
Oznaczmy $a_n=(A_n-A'_n)/2$, zatem $a_n\to0$
w $W^{k,2}$. Ponieważ $|u_n|=1$,
normy $\|u_n\|_{L^2}$ są stałe.
Ciągłość mnożenia przez formę klasy $W^{k,2}$
na $W^{j,2}$ dla $0\le j\le k$ oraz
$\mathrm du_n=u_na_n$ dają indukcyjnie
jednostajne ograniczenia $u_n$ w $W^{j+1,2}$
i $\mathrm du_n\to0$ w $W^{j,2}$.
W szczególności nierówność Poincarégo dla funkcji
o średniej zero, zastosowana do obu składowych,
daje $u_n-\overline u_n\to0$ w $W^{k+1,2}$,
gdzie $\overline u_n$ oznacza średnią całkową
funkcji $u_n$, a nie jej sprzężenie zespolone.
Po wydzieleniu podciągu średnie zbiegają
do stałej $u_\infty\in S^1$ (z warunku $|u_n|=1$
i zbieżności w $C^0$).
Przechodząc do granicy w równaniu spinorowym,
mamy $u_\infty\psi=\psi$, więc $u_\infty=1$.
Zatem dla dostatecznie bliskich konfiguracji
każda transformacja, która je utożsamia,
należy do ustalonego otoczenia identyczności;
w przeciwnym razie powyższy argument dałby
sprzeczność na ciągu malejących otoczeń.

Po zmniejszeniu otoczenia przekroju twierdzenie
o funkcji uwikłanej zapewnia więc, że dwie jego
konfiguracje leżące na tej samej orbicie są równe.
Wraz z gładką lokalną funkcją $f(\xi)$ daje to
lokalny model gładkiego ilorazu. Uzasadnijmy jeszcze
własność Hausdorffa: jeśli $\mathfrak c_n\to\mathfrak c$
i $\mathfrak c'_n\to\mathfrak c'$ oraz
$u_n\cdot\mathfrak c_n=\mathfrak c'_n$, to tym razem
$a_n=(A_n-A'_n)/2$ zbiega w $W^{k,2}$, niekoniecznie
do zera. Ta sama indukcja daje ograniczenie $u_n$
w $W^{k+1,2}$. Z twierdzenia Rellicha wydzielamy
podciąg zbieżny w $W^{k,2}$ do $u$.
Zanurzenie w $C^0$ daje $|u|=1$, a ciągłość
mnożenia w $W^{k,2}$ daje zbieżność
$\mathrm du_n=u_na_n$ w $W^{k,2}$.
W konsekwencji $u_n\to u$ w $W^{k+1,2}$,
więc $\mathfrak c'=u\cdot\mathfrak c$.
Relacja równoważności jest zatem domknięta
w iloczynie przestrzeni konfiguracji.
Ponieważ rzut ilorazowy jest otwarty, iloraz
jest przestrzenią Hausdorffa i ma przeliczalną
bazę jako otwarty obraz ośrodkowej przestrzeni
Hilberta (afinicznej). Na części nieredukowalnej
powyższe przekroje dają mapy gładkiej rozmaitości;
zmiany współrzędnych są gładkie przez lokalne
twierdzenie o funkcji uwikłanej, po przesunięciu
reprezentantów ustaloną transformacją cechowania.
Nie oznacza to istnienia jednego przekroju globalnego.
\end{proof}

\begin{stwierdzenie}[Krata okresów grupy cechowania]
\label{prop:tg-sw12-lattice}
Niech $u\in C^\infty(X,S^1)$. Forma
$\alpha_u:=\frac{1}{2\pi i}u^{-1}\mathrm du$ jest zamknięta
i reprezentuje klasę w obrazie
$H^1(X;\mathbb Z)\to H^1(X;\mathbb R)$.
Dla $A=A_0+i a$, $a\in\Omega^1(X)$, transformacja $u$
zmienia część rzeczywistą $a$ o $-4\pi\alpha_u$.
W konsekwencji harmoniczna część $a$ jest określona
modulo kratę $4\pi H^1(X;\mathbb Z)$, rozumianą
w przestrzeni harmonicznych form stopnia jeden.
\end{stwierdzenie}

\begin{proof}
Ograniczenie funkcji $u$ do dowolnej gładkiej pętli
$\gamma$ ma całkowity stopień, więc
$\int_\gamma u^{-1}\mathrm du/(2\pi i)\in\mathbb Z$.
Klasa formy jest tym samym całkowita. Ze wzoru
$u\cdot A=A-2u^{-1}\mathrm du$ otrzymujemy
$u\cdot A=A_0+i(a-4\pi\alpha_u)$.
Odwrotnie, niech $\alpha$ będzie harmonicznym reprezentantem
klasy całkowitej. Ustalmy $x_0\in X$ i połóżmy
$u(x)=\exp(2\pi i\int_{x_0}^x\alpha)$,
gdzie całkujemy wzdłuż dowolnej odcinkami gładkiej
drogi. Niezależność od drogi wynika z całkowitości
okresów; lokalnie $\alpha=\mathrm df$, więc
$u$ jest gładka i $u^{-1}\mathrm du=2\pi i\alpha$.
Każde przesunięcie przez harmoniczny reprezentant
kraty $4\pi H^1(X;\mathbb Z)$ realizuje zatem
transformacja cechowania zachowująca warunek Coulomba.
Jeżeli dwie transformacje dają tę samą klasę,
różnica ich form ma zerowe okresy i jest dokładna.
To daje dokładnie wskazaną kratę.
\end{proof}

\begin{przyklad}[Cechowanie harmoniczne na $S^1\times S^3$]
Niech $t\in\mathbb R/\mathbb Z$ będzie współrzędną na
$S^1$ i niech $X=S^1\times S^3$ z metryką produktową.
Dla $u_n(t,x)=e^{2\pi int}$, $n\in\mathbb Z$,
otrzymujemy $u_n^{-1}\mathrm du_n=2\pi in\,\mathrm dt$.
Połączenia $A_0+i\theta\,\mathrm dt$ i
$A_0+i(\theta-4\pi n)\,\mathrm dt$ są cechowaniowo
równoważne. Parametr harmoniczny należy więc do okręgu
$\mathbb R/(4\pi\mathbb Z)$, a nie do całej prostej.
Ten przykład dotyczy mechanizmu cechowania; sam w sobie
nie wyznacza żadnego niezmiennika Seiberga--Wittena.
\end{przyklad}

\begin{center}
\begin{tikzpicture}[>=Stealth,scale=.9]
  \fill[tgBlue!10] (-4.2,-1.3) rectangle (4.2,1.5);
  \draw[tgNavy!50,->] (-4.45,-.05)--(4.5,-.05);
  \draw[tgPurple!75,dashed,line width=1.3pt]
    (-3.55,-1.28)--(-3.55,1.45)
    (3.55,-1.28)--(3.55,1.45);
  \draw[tgTeal,line width=2.6pt] (-3.55,.48)--(3.55,.48);
  \draw[tgCoral,-{Stealth[length=2.8mm]},line width=2pt]
    (3.05,.72) .. controls (4.1,1.18) and (4.1,-.65) .. (3.18,-.84);
  \node[font=\small,text=tgNavy] at (0,1.06)
    {parametr harmoniczny $h$};
  \node[font=\scriptsize,text=tgPurple] at (-3.55,-1.52) {$h$};
  \node[font=\scriptsize,text=tgPurple] at (3.55,-1.52) {$h+4\pi\alpha$};
  \node[font=\scriptsize,text=tgCoral] at (2.0,-1.0)
    {utożsamienie przez cechowanie};
\end{tikzpicture}

{\footnotesize Schemat kraty okresów w przestrzeni części harmonicznych
połączenia. Pokazano jeden kierunek kraty; przy $b_1(X)>1$
otrzymujemy torus wyższego wymiaru.}
\end{center}

\section{Symbol kompleksu i indeks bez pomijania przejść}
\label{sec:tg-sw12-symbol}

\begin{lemat}[Dokładność symbolu kompleksu samodualnego]
\label{lem:tg-sw12-symbol}
Niech $V$ będzie zorientowaną euklidesową przestrzenią
rzeczywistą wymiaru cztery, a $\xi\in V^*\setminus\{0\}$.
Wówczas ciąg
\[
0\longrightarrow\mathbb R
\xrightarrow{f\mapsto f\xi}V^*
\xrightarrow{a\mapsto(\xi\wedge a)^+}\Lambda^2_+V^*
\longrightarrow0
\]
jest dokładny.
\end{lemat}

\begin{proof}
Normalizując $\xi$, możemy przyjąć $\xi=e^1$ w dodatnio
zorientowanej bazie ortonormalnej $e^1,\dots,e^4$.
Dla $a=\sum_{j=1}^4a_je^j$ dostajemy
\[
(e^1\wedge a)^+=\frac12\bigl[
 a_2(e^{12}+e^{34})+a_3(e^{13}-e^{24})
 +a_4(e^{14}+e^{23})\bigr].
\]
Trzy formy w nawiasach stanowią bazę $\Lambda^2_+V^*$.
Odwzorowanie drugie jest więc surjektywne i ma jądro
$\mathbb Re^1$, dokładnie równe obrazowi pierwszego.
\end{proof}

\begin{lemat}[Jawna macierz symbolu operatora samodualnego]
\label{lem:tg-sw12-symbol-matrix}
Niech $V$ będzie zorientowaną rzeczywistą przestrzenią
euklidesową wymiaru cztery, z dodatnio zorientowaną
bazą ortonormalną $(e^1,e^2,e^3,e^4)$ przestrzeni $V^*$.
Niech $\xi=\sum_{j=1}^4\xi_je^j$ i użyjmy bazy
$(e^{12}+e^{34},\,e^{13}-e^{24},\,e^{14}+e^{23})$
przestrzeni $\Lambda^2_+V^*$.
Macierz odwzorowania
\[
V^*\longrightarrow\mathbb R\oplus\Lambda^2_+V^*,
\qquad a\longmapsto
(-\langle\xi,a\rangle,(\xi\wedge a)^+)
\]
w podanych bazach jest równa
\[
M(\xi)=\begin{pmatrix}
-\xi_1&-\xi_2&-\xi_3&-\xi_4\\
-\xi_2/2&\xi_1/2&-\xi_4/2&\xi_3/2\\
-\xi_3/2&\xi_4/2&\xi_1/2&-\xi_2/2\\
-\xi_4/2&-\xi_3/2&\xi_2/2&\xi_1/2
\end{pmatrix}.
\]
Dla każdego $\xi\in V^*$ zachodzi
\[
M(\xi)M(\xi)^T
=|\xi|^2\operatorname{diag}(1,\tfrac14,\tfrac14,\tfrac14).
\]
Stąd dla $\xi\ne0$ symbol jest odwracalny,
a $|\det M(\xi)|=|\xi|^4/8$.
\end{lemat}

\begin{proof}
Część zerowego stopnia daje pierwszy wiersz macierzy.
Rzut na przestrzeń form samodualnych zapisujemy jako
$\alpha^+=(\alpha+*\alpha)/2$; obliczenie współczynników
w podanej bazie daje pozostałe trzy wiersze.
Pierwszy wiersz ma kwadrat normy $|\xi|^2$,
a każdy z trzech pozostałych ma kwadrat normy
$|\xi|^2/4$. Są one parami ortogonalne:
na przykład iloczyn pierwszego z drugim wynosi
$(\xi_1\xi_2-\xi_2\xi_1+
\xi_3\xi_4-\xi_4\xi_3)/2=0$.
Pozostałe iloczyny znikają analogicznie,
co dowodzi tożsamości macierzowej.
Dla $\xi\ne0$ macierz iloczynu
$M(\xi)M(\xi)^T$ jest dodatnio określona,
więc $M(\xi)$ ma pełny rząd.
Biorąc wyznaczniki, otrzymujemy
$(\det M(\xi))^2=|\xi|^8/64$,
skąd żądany wzór.
Rachunek ten rozwija Lemat~\ref{lem:tg-sw12-symbol}
do tożsamości ilościowej, bez użycia
twierdzeń analitycznych o operatorach Fredholma.

Należy rozróżnić normę współrzędnych od normy
form samodualnych: trzy wektory użytej bazy
mają normę $\sqrt2$, a nie jeden.
W bazie ortonormalnej przestrzeni docelowej
macierz symbolu jest zatem równa
$N(\xi)=\operatorname{diag}(1,\sqrt2,\sqrt2,\sqrt2)M(\xi)$.
Stąd
\[
N(\xi)N(\xi)^T=|\xi|^2
\operatorname{diag}(1,\tfrac12,\tfrac12,\tfrac12).
\]
Dla $\xi\ne0$ wartości własne $N(\xi)^TN(\xi)$
są takie same jak wartości własne $N(\xi)N(\xi)^T$,
gdyż macierze te są podobne przez $N(\xi)$.
Otrzymujemy więc oszacowania w naturalnych normach
\[
\tfrac12|\xi|^2|a|^2
\leq |\langle\xi,a\rangle|^2+|(\xi\wedge a)^+|^2
\leq |\xi|^2|a|^2.
\]
Dla $\xi=0$ obie strony znikają. Stałe są ostre:
dla $\xi\ne0$ górna odpowiada $a$ równoległemu
do $\xi$, a dolna niezerowemu $a$ ortogonalnemu
do $\xi$, jak w obliczeniu
z lematu \ref{lem:tg-sw12-symbol}.
\end{proof}

\begin{twierdzenie}[Eliptyczność operatora cechowania i równań]
\label{thm:tg-sw12-fredholm}
Dla rozwiązania $(A,\psi)$ równań Seiberga--Wittena operator
\[
\mathcal D_{A,\psi}:=\mathrm d_0^*\oplus\mathrm d_1:
\Omega^1(X;i\mathbb R)\oplus\Gamma(S^+)
\longrightarrow
\Omega^0(X;\mathbb R)\oplus\Gamma(S^-)
\oplus\Omega^2_+(X;i\mathbb R)
\]
jest eliptyczny. Po dopełnieniu w przestrzeniach Sobolewa
jest operatorem Fredholma o indeksie
$\operatorname{ind}_\mathbb R\mathcal D_{A,\psi}
=\bigl(c_1(L)^2-2e(X)-3\sigma(X)\bigr)/4$.
\end{twierdzenie}

\begin{proof}
Człony zależne od $\psi$ mają rząd zero, więc symbol
operatora, po przestawieniu czynników i
utożsamieniu form o wartościach w $i\mathbb R$
z formami rzeczywistymi, jest sumą prostą
symboli $(-2\mathrm d^*)\oplus\mathrm d^+$
oraz Diraca. Czynnik $-2$ w pierwszym
składniku jest odwracalny, więc nie wpływa
na eliptyczność ani na indeks.
Lemat \ref{lem:tg-sw12-symbol} pokazuje, że pierwszy jest
izomorfizmem dla każdego niezerowego kowektora: jego jądro
ma wymiar zero, a dziedzina i przeciwdziedzina mają
wymiar cztery. Drugi symbol jest ograniczeniem mnożenia Clifforda
$\rho(\xi)\colon S^+\to S^-$. Jest odwracalny,
ponieważ $\rho(\xi)^*\rho(\xi)=
|\xi|^2\operatorname{id}_{S^+}$ dla $\xi\ne0$.
Wynika to z relacji Clifforda na całej wiązce spinorowej;
obie półwiązki mają ten sam wymiar zespolony.
Dostajemy eliptyczność.

Po odkształceniu członów rzędu zero indeks nie zmienia
się. Indeks rzeczywisty części $\mathrm d^*\oplus\mathrm d^+$
wynosi $b_1(X)-b_0(X)-b_2^+(X)=b_1-1-b_2^+$.
Indeks rzeczywisty Diraca jest dwukrotnością indeksu
zespolonego, a twierdzenie Atiyaha--Singera daje
$2\operatorname{ind}_\mathbb C D_A=(c_1(L)^2-\sigma)/4$.
Po dodaniu i użyciu $e=2-2b_1+b_2^++b_2^-$,
$\sigma=b_2^+-b_2^-$, otrzymujemy wzór z tezy.
Twierdzenie indeksowe jest tu wyraźnie przyjęte z literatury.
\end{proof}

\begin{przyklad}[Rachunek wymiaru na rozdmuchaniach $\mathbb{CP}^2$]
Niech $X_k=\mathbb{CP}^2\#k\overline{\mathbb{CP}}^{\,2}$
i $c=aH-\sum_{j=1}^k b_jE_j$, gdzie wszystkie liczby
$a,b_j$ są nieparzyste. Wtedy
\[
d(c)=\frac{a^2-\sum_{j=1}^k b_j^2-9+k}{4}.
\]
Dla $a=3$ i $b_j=1$ otrzymujemy $d(c)=0$ dla każdego $k$.
Równość wymiaru oczekiwanego zeru nie mówi jednak,
że istnieje rozwiązanie: przy metryce o dodatniej
krzywiźnie skalarnej i $\eta=0$ przestrzeń może być pusta.
\end{przyklad}

\section{Zwartość: od ograniczenia spinora do zbieżności konfiguracji}
\label{sec:tg-sw12-compact}

Największy skok w standardowych krótkich dowodach
zwartości polega na przejściu od ograniczenia $|\psi|$
do zwartości \emph{połączeń}. Rozpiszemy go przy użyciu
klasycznej teorii Hodge'a i oszacowań eliptycznych.

\begin{lemat}[Kontrola pełnej krzywizny]
\label{lem:tg-sw12-curvature}
Dla połączenia unitarnego $A$ na $L$ zachodzi
\[
\|F_A\|_{L^2}^2=2\|F_A^+\|_{L^2}^2
-4\pi^2\langle c_1(L)^2,[X]\rangle.
\]
W szczególności przy ustalonej klasie $c_1(L)$
i jednostajnie ograniczonej normie $\|F_A^+\|_{L^2}$
pełna krzywizna jest jednostajnie ograniczona w $L^2$.
\end{lemat}

\begin{proof}
Piszemy $F_A=i\Omega$, gdzie $\Omega$ jest rzeczywistą
zamkniętą formą. Z teorii Cherna--Weila
$[-\Omega/(2\pi)]=c_1(L)\otimes1$
w konwencji z definicji \ref{def:tg-sw-eqn}.
Znak znika przy obliczeniu kwadratu klasy.
Rozkład Hodge'a daje
\[
\int_X\Omega\wedge\Omega
=\|\Omega^+\|_{L^2}^2-\|\Omega^-\|_{L^2}^2
=4\pi^2\langle c_1(L)^2,[X]\rangle.
\]
Dodając i odejmując kwadraty obu norm, otrzymujemy wzór.
\end{proof}

\begin{twierdzenie}[Zwartość rozwiązań przy ustalonych danych]
\label{thm:tg-sw12-compact}
Niech $g$, $\mathfrak s$ i gładka forma $\eta$ będą ustalone.
Dla dowolnego ciągu gładkich rozwiązań
$(A_j,\psi_j)$ równań Seiberga--Wittena istnieją podciąg
i gładkie transformacje cechowania $u_j$ takie, że
$u_j\cdot(A_j,\psi_j)$ zbiega w topologii $C^\infty$
do gładkiego rozwiązania. W szczególności przestrzeń
moduli $\mathcal M_{\mathfrak s}(g,\eta)$ jest zwarta.
\end{twierdzenie}

\begin{proof}[Pełny schemat z jawnym użyciem oszacowań eliptycznych]
\emph{Krok 1.} Z twierdzenia \ref{thm:tg-sw12-spinor-bound}
spinory $\psi_j$ są jednostajnie ograniczone w $L^\infty$.
Równanie $F_{A_j}^+=q(\psi_j)+i\eta$ daje zatem
jednostajne ograniczenie $F_{A_j}^+$ w $L^\infty$.
Lemat \ref{lem:tg-sw12-curvature} zapewnia dodatkowo
ograniczenie pełnej krzywizny w $L^2$.

\emph{Krok 2.} Wybieramy gładkie połączenie $A_0$ na $L$
i zapisujemy $A_j=A_0+i a_j$.
Rozkład Hodge'a formy rzeczywistej $a_j$ ma postać
$a_j=\mathrm df_j+\mathrm d^*\beta_j+h_j$.
Transformacja $e^{if_j/2}$ usuwa część dokładną.
Następnie transformacje odpowiadające klasom całkowitym
pozwalają wybrać części harmoniczne $h_j$ w ustalonym
domkniętym równoległościanie fundamentalnym kraty ze
stwierdzenia \ref{prop:tg-sw12-lattice}. Otrzymujemy
$\mathrm d^*a_j=0$ i jednostajnie ograniczone $h_j$.

\emph{Krok 3.} Z równania krzywizny mamy
$\mathrm d^+a_j=(F_{A_j}^+-F_{A_0}^+)/i$.
Jego prawa strona jest ograniczona w każdym $L^p$,
$1<p<\infty$. Eliptyczne oszacowanie dla operatora
$\mathrm d^*\oplus\mathrm d^+$ daje
\[
\|a_j\|_{W^{1,p}}\le C_p
\bigl(\|\mathrm d^+a_j\|_{L^p}
+\|\mathrm d^*a_j\|_{L^p}+\|h_j\|_{L^p}\bigr)\le C'_p.
\]
Wybieramy $p>4$ oraz $0<\alpha<1-4/p$.
Z zanurzenia Sobolewa
$W^{1,p}\hookrightarrow C^{0,\alpha}$ otrzymujemy
jednostajny moduł ciągłości form $a_j$.

\emph{Krok 4.} Względem połączenia odniesienia
$D_{A_0}\psi_j=-\tfrac12\rho(i a_j)\psi_j$.
Prawa strona jest ograniczona w $L^p$ dzięki
$\|a_j\|_{L^p}$ i $\|\psi_j\|_{L^\infty}$.
Z oszacowania eliptycznego
\[
\|\psi_j\|_{W^{1,p}}\leq C_p
(\|D_{A_0}\psi_j\|_{L^p}+\|\psi_j\|_{L^p})
\]
wynika jednostajne ograniczenie $\psi_j$ w $W^{1,p}$;
ostatnia norma po prawej stronie jest ograniczona
dzięki oszacowaniu w $L^\infty$ i skończonej objętości $X$.
Ponieważ $q$ jest wielomianem kwadratowym, z tego
ograniczenia wynika $q(\psi_j)$ ograniczone w $W^{1,p}$.
Stąd prawa strona równania dla $\mathrm d^+a_j$ jest
ograniczona w $W^{1,p}$ i eliptyczność
$\mathrm d^*\oplus\mathrm d^+$ daje $a_j$ ograniczone w
$W^{2,p}$. Dokładniej, dla każdego $m\geq1$ przestrzeń $W^{m,p}$
jest algebrą, gdyż $p>4$. Jeżeli $a_j,\psi_j$
są ograniczone w $W^{m,p}$, to obie prawe strony
$\mathrm d^+a_j=(q(\psi_j)+i\eta-F_{A_0}^+)/i$
oraz $D_{A_0}\psi_j=-\rho(ia_j)\psi_j/2$
są ograniczone w $W^{m,p}$.
Oszacowania eliptyczne rzędu $m+1$ dla obu
operatorów, z dodatkowymi ograniczonymi normami
$L^p$ lub częściami harmonicznymi, dają
ograniczenia obu ciągów w $W^{m+1,p}$.
Indukcja zaczyna się od $m=1$ i nie wymaga
zmiany cechowania w kolejnych krokach.

\emph{Krok 5.} Z twierdzenia Rellicha i osadzeń Sobolewa
wydzielamy metodą przekątniową podciąg zbieżny w $C^m$
dla każdego $m$. Granica spełnia równania, ponieważ
operator Diraca, krzywizna i odwzorowanie $q$ są ciągłe
w tych topologiach. Przejście do wszystkich klas
rozwiązań Sobolewa i od zwartości ciągowej
do zwartości ilorazu uzasadniono w dowodzie
twierdzenia \ref{thm:tg-sw-compact}.
W ten sposób otrzymujemy zwartość modulo cechowanie.

W krokach 3--5 korzystaliśmy z \emph{standardowych globalnych
oszacowań eliptycznych} i twierdzeń Sobolewa--Rellicha.
Ich pełne dowody nie należą do teorii grafów ani do
algebry spinorowej; argument wskazuje natomiast dokładnie,
które szacunki są potrzebne i dlaczego nie wystarczy samo
ograniczenie spinora.
\end{proof}

\begin{uwaga}[Dlaczego nie pojawia się koncentracja energii?]
W teorii połączeń o nieabelowej grupie strukturalnej
zbieżność może wymagać usunięcia punktów koncentracji.
Tutaj mamy grupę $U(1)$, równanie samodualnej krzywizny
bez wyrazu $A\wedge A$ oraz bezpośredni szacunek
$F_A^+=q(\psi)+i\eta$. Po ustaleniu cechowania jest to
równanie eliptyczne dla formy $a$, a jego oszacowania
dają zbieżność bez dodawania idealnych połączeń.
\end{uwaga}

\section{Transwersalność dla nieredukowalnych rozwiązań}
\label{sec:tg-sw12-transversality}

Zwartość nie przesądza, że przestrzeń moduli jest gładka.
Aby liczyć rozwiązania ze znakami, trzeba jeszcze wykazać
surjektywność linearyzacji, zwykle po generycznym wyborze
perturbacji.

\begin{lemat}[Niezdegenerowanie mnożenia Clifforda przez spinor]
\label{lem:tg-sw12-clifford-real-iso}
Niech $x\in X$, $\psi\in S_x^+$ będzie
niezerowym spinorem, a
$\rho:T_x^*X\otimes_{\mathbb R}\mathbb C
\to\operatorname{Hom}_{\mathbb C}(S_x^+,S_x^-)$
oznacza mnożenie Clifforda.
Wówczas odwzorowanie rzeczywisto-liniowe
\[
T_x^*X\longrightarrow S_x^-,
\qquad v\longmapsto\rho(iv)\psi
\]
jest izomorfizmem rzeczywistych przestrzeni
wektorowych wymiaru cztery.
\end{lemat}

\begin{proof}
Relacje Clifforda dla rzeczywistego
kowektora $v$ zapewniają
$\rho(v)^*\rho(v)=|v|^2\operatorname{id}_{S_x^+}$
(przy konwencji hermitowskiego sprzężenia).
Ponieważ $\rho(iv)=i\rho(v)$
i mnożenie przez $i$ nie zmienia normy,
mamy
\[
|\rho(iv)\psi|^2
=\langle\rho(iv)^*\rho(iv)\psi,\psi\rangle
=|v|^2|\psi|^2.
\]
Jeżeli $\rho(iv)\psi=0$, to
$|v|^2|\psi|^2=0$. Z założenia
$\psi\ne0$, więc $v=0$;
odwzorowanie jest zatem injektywne.
Ponieważ $\dim_{\mathbb R}T_x^*X=4$
i $\dim_{\mathbb R}S_x^-=4$,
jest również surjektywne.
Dowód wykorzystuje wyłącznie
lokalną algebrę Clifforda, a nie
twierdzenie o jednoznacznym przedłużeniu.
\end{proof}

\begin{lemat}[Cokernel uniwersalnej linearyzacji]
\label{lem:tg-sw12-surj}
Niech $(A,\psi)$ będzie gładkim rozwiązaniem nieredukowalnym.
Rozważmy linearyzację operatora równań, w której oprócz
zmiany $(A,\psi)$ dopuszczamy dowolną zmianę perturbacji
$\eta\in\Omega^2_+(X)$. Wtedy operator
\[
(a,\phi,\zeta)\longmapsto
\bigl(D_A\phi+\tfrac12\rho(a)\psi,\,
\mathrm d^+a-(\mathrm dq)_\psi\phi-i\zeta\bigr)
\]
ma gęsty obraz w rzeczywistej przestrzeni $L^2$
przekrojów $S^-\oplus i\Lambda^2_+$.
Dla $k\geq3$, po ograniczeniu wariacji konfiguracji
do $\ker\mathrm d_0^*$ klasy $W^{k,2}$ oraz przy
perturbacjach klasy $W^{k-1,2}(\Lambda^2_+)$,
operator uniwersalny jest surjektywny na
$W^{k-1,2}(S^-\oplus i\Lambda^2_+)$.
\end{lemat}

\begin{proof}
Przypuśćmy, że para $(\chi,\beta)$ jest ortogonalna
do obrazu. Testując dowolne $\zeta\in\Omega^2_+(X)$,
otrzymujemy $\beta=0$. Testując dowolne $\phi$,
otrzymujemy $D_A^*\chi=0$ w sensie dystrybucyjnym,
a z regularności eliptycznej $\chi$ jest gładka.
Testując wszystkie $a\in\Omega^1(X;i\mathbb R)$,
dostajemy punktowo
$\operatorname{Re}\langle\chi,\rho(a)\psi\rangle=0$
dla wszystkich $a$. W każdym punkcie, w którym $\psi\ne0$,
odwzorowanie rzeczywiste
$T_x^*X\to S_x^-$ dane przez $v\mapsto\rho(iv)\psi$
jest izomorfizmem na mocy Lematu~\ref{lem:tg-sw12-clifford-real-iso}.
Stąd $\chi=0$ na niepustym zbiorze otwartym
$\{\psi\ne0\}$. Dla operatora Diraca obowiązuje
zasada jednoznacznego przedłużenia: gładki spinor w jądrze
$D_A^*$, znikający na niepustym zbiorze otwartym,
znika na spójnej rozmaitości. Wynika zatem $\chi=0$.

Otrzymaliśmy brak niezerowej pary ortogonalnej do obrazu.
Aby przejść od gęstości w $L^2$ do surjektywności
w normie Sobolewa, nie wystarcza sama gęstość.
Niech $T=\mathrm d_1|_{\ker\mathrm d_0^*}$.
Odwracalność $\mathrm d_0^*\mathrm d_0$ w punkcie
nieredukowalnym (lemat \ref{lem:tg-sw-full-coulomb})
daje rozkład każdej wariacji na element
$\ker\mathrm d_0^*$ oraz element
$\operatorname{im}\mathrm d_0$.
Ponieważ $\mathrm d_1\mathrm d_0=0$, mamy
$\operatorname{im}T=\operatorname{im}\mathrm d_1$.
Eliptyczność kompleksu z twierdzenia
\ref{thm:tg-sw-elliptic} zapewnia domkniętość
tego obrazu i skończony wymiar cokernelu.
Nie korzystamy tutaj z twierdzenia
o generycznej regularności.
Eliptyczny rozkład Hodge'a dla kompleksu
deformacyjnego reprezentuje ten cokernel
przez gładkie przekroje w $\ker\mathrm d_1^*$.
Dokładniej, dla gładkiego rozwiązania operator
$\mathrm d_1\mathrm d_1^*$ na ostatnim wyrazie
kompleksu jest eliptyczny, a jego harmoniczne
jądro daje reprezentanty cokernelu także
w wyższych klasach Sobolewa, przez regularność
eliptyczną. Jeżeli $P\zeta=(0,-i\zeta)$,
to odwzorowanie $\zeta\mapsto[P\zeta]$
na ten skończeniewymiarowy cokernel
jest surjektywne: w przeciwnym razie
istniałby niezerowy gładki reprezentant
ortogonalny zarówno do $\operatorname{im}\mathrm d_1$,
jak i do wszystkich $P\zeta$.
Wyklucza to pierwsza część dowodu.
Dla dowolnego prawego członu $y$ wybieramy
więc $\zeta$ z $[P\zeta]=[y]$.
Wówczas $y-P\zeta\in\operatorname{im}T$,
co daje żądaną surjektywność.
\end{proof}

\begin{twierdzenie}[Generyczna regularność; użycie Sarda--Smale'a]
\label{thm:tg-sw12-regular}
Załóżmy $b_2^+(X)>0$ i ustalmy metrykę $g$ oraz
strukturę $\mathfrak s$. Dla generycznej gładkiej
samodualnej perturbacji $\eta$ przestrzeń moduli
$\mathcal M_{\mathfrak s}(g,\eta)$ nie ma rozwiązań
redukowalnych, a wszystkie jej rozwiązania są regularne.
Jeśli jest niepusta, jest więc zwartą, gładką
zorientowalną rozmaitością rzeczywistą wymiaru
$d(\mathfrak s)$.
\end{twierdzenie}

\begin{proof}[Redukcja do twierdzeń analitycznych]
Z równania ściany (stwierdzenie \ref{prop:tg-sw-wall})
wynika, że perturbacje dopuszczające rozwiązanie
redukowalne tworzą afiniczną podprzestrzeń kowymiaru
$b_2^+(X)>0$. Poza nią wszystkie rozwiązania są
nieredukowalne. Lemat \ref{lem:tg-sw12-surj} zapewnia
surjektywność uniwersalnej linearyzacji przy takich
rozwiązaniach. W otoczeniu każdego gładkiego rozwiązania
uniwersalna linearyzacja na dopełnieniach Sobolewa
jest surjektywna i ma dopełnione jądro.
Twierdzenie o funkcji uwikłanej i lokalny przekrój
dają tam gładki lokalny model przestrzeni uniwersalnej.
Różniczka rzutowania na perturbacje ma indeks
$d(\mathfrak s)$, jak wykazano w dowodzie twierdzenia
\ref{thm:tg-sw-generic}.
Nie potrzebujemy tu twierdzenia o gładkości
przestrzeni uniwersalnej dla wszystkich perturbacji
o ograniczonej regularności Sobolewa.
Aby otrzymać generyczność w topologii $C^\infty$,
możemy posłużyć się następującym argumentem
ze skończeniewymiarowymi perturbacjami.
Ustalmy dowolną gładką perturbację $\eta_0$
poza ścianą. Dla każdej klasy rozwiązania
przy $\eta_0$ wybieramy skończenie wiele
gładkich form samodualnych, których obrazy
rozpinają cokernel linearyzacji w tym punkcie.
Jest to możliwe na mocy dowodu lematu
\ref{lem:tg-sw12-surj}, który używa właśnie
gładkich wariacji perturbacji.
Surjektywność rozszerzonej linearyzacji
jest warunkiem otwartym: surjektywny operator
między przestrzeniami Hilberta ma ograniczoną
prawą odwrotność, a dostatecznie mała zmiana
operatora zachowuje jej istnienie przez szereg
Neumanna. W lokalnych przekrojach otrzymujemy
więc otoczenia rozwiązań, na których te same
skończone rodziny kierunków wystarczają.
Zwartość $\mathcal M_{\mathfrak s}(g,\eta_0)$
pozwala wybrać skończone pokrycie tymi otoczeniami.
Niech $P$ będzie liniową rozpiętością wszystkich
wybranych form, a dla $h\in P$ połóżmy
$\eta_h=\eta_0+h$.

Dla dostatecznie małego otoczenia zera w $P$
wszystkie rozwiązania przy $\eta_h$ są objęte
tymi lokalnymi modelami i rozszerzona
linearyzacja jest surjektywna.
W przeciwnym razie istniałby ciąg $h_n\to0$
i rozwiązań poza tymi otoczeniami.
Ponieważ $P$ ma skończony wymiar i składa się
z form gładkich, $\eta_{h_n}\to\eta_0$
w każdej normie $C^m$. Dowód zwartości
z twierdzenia \ref{thm:tg-sw12-compact}
daje wtedy po cechowaniu podciąg zbieżny
do rozwiązania przy $\eta_0$: wszystkie
użyte tam oszacowania są jednostajne dla
takiej rodziny perturbacji.
Byłaby to sprzeczność z wyborem pokrycia.
Ten argument obejmuje także pustą przestrzeń
moduli przy $\eta_0$, która pozostaje pusta
dla dostatecznie bliskich gładkich danych.

Uniwersalna przestrzeń rozwiązań nad małym
otoczeniem zera w $P$ jest zatem, jeśli niepusta,
gładką rozmaitością skończonego wymiaru
$d(\mathfrak s)+\dim P$. Wszystkie rozwiązania
mają gładkich reprezentantów, bo dane są gładkie.
Wartość regularna rzutowania na $P$ odpowiada
surjektywności linearyzacji przy ustalonej
perturbacji: surjektywność rozszerzonego operatora
zapewnia surjektywność mapy z $P$ na jego
nierozszerzony cokernel, a obraz różniczki
rzutowania jest jądrem tej ostatniej mapy.
Z twierdzenia Sarda dla jej rzutowania na $P$
otrzymujemy gęsty zbiór wartości regularnych.
Możemy wybrać $h$ dowolnie bliskie zeru
w $P$, a tym samym w dowolnym zadanym
otoczeniu zera w topologii $C^\infty$.
Regularność wszystkich rozwiązań jest
również otwarta poza ścianą w tej topologii:
gdyby ciąg nieregularnych rozwiązań dla
$\eta_n\to\eta$ zbiegał po cechowaniu
do regularnego rozwiązania dla $\eta$,
otwartość surjektywności linearyzacji
dałaby sprzeczność. Istnienie takiego
podciągu zapewnia ponownie powyższy
argument zwartości. Otrzymujemy więc
otwarty i gęsty, w szczególności rezydualny,
zbiór gładkich perturbacji regularnych.
Jest to uzasadnienie przejścia do danych
gładkich, niezależne od samej gęstości
gładkich form w dopełnieniu Sobolewa.
Twierdzenie \ref{thm:tg-sw12-compact} zapewnia zwartość.
Orientowalność wynika z opisu linii wyznacznikowej
w kolejnej sekcji. W nieskończenie wymiarowej rodzinie perturbacji
odpowiednim narzędziem jest twierdzenie
Sarda--Smale'a; powyższa redukcja dla ustalonej
metryki i lokalnej rodziny skończeniewymiarowej
wykorzystuje zwykłe twierdzenie Sarda.
\end{proof}

\section{Orientacja przestrzeni moduli i zliczanie rozwiązań}
\label{sec:tg-sw12-orientation}

\begin{definicja}[Linia wyznacznikowa operatora Fredholma]
Niech $D:E\to F$ będzie rzeczywistym operatorem Fredholma
między przestrzeniami Banacha. Jego \emph{linią wyznacznikową}
nazywamy jednowymiarową rzeczywistą przestrzeń
$\det D:=\Lambda^{\max}\ker D\otimes
(\Lambda^{\max}\operatorname{coker}D)^*$.
Przez $\Lambda^{\max}0$ rozumiemy $\mathbb R$.
\end{definicja}

\begin{stwierdzenie}[Wybór orientacji homologicznej]
\label{prop:tg-sw12-hom-orient}
Po wybraniu orientacji rzeczywistej przestrzeni
$H^1(X;\mathbb R)\oplus H^2_+(X;\mathbb R)$
wszystkie regularne nieredukowalne przestrzenie moduli
Seiberga--Wittena na $X$ otrzymują zgodne orientacje.
\end{stwierdzenie}

\begin{proof}[Opis linii orientacji]
Operator deformacyjny $\mathcal D_{A,\psi}$ można odkształcić
przez człony rzędu zero do sumy
$((-2\mathrm d^*)\oplus\mathrm d^+)\oplus D_A$,
po identyfikacji form imaginarnych z rzeczywistymi
i ustaleniu kolejności składników.
Nie usuwamy znaku $-2$ przez homotopię skalarów
przechodzącą przez zero: identyfikujemy
odpowiednie jądra i cokernele bezpośrednio.
Wyznacznik zespolonego operatora $D_A$ ma orientację
kanoniczną, ponieważ jądro i cokernel są zespolonymi
przestrzeniami wektorowymi. Linia wyznacznikowa
$\mathrm d^*\oplus\mathrm d^+$ odpowiada, przez teorię
Hodge'a, przestrzeniom
$H^1(X;\mathbb R)$, $H^0(X;\mathbb R)$ i $H^2_+(X;\mathbb R)$.
Czynnik $H^0$ ma orientację kanoniczną wyznaczoną przez
funkcję stałą $1$, a orientacja pozostałych dwóch
przestrzeni wyznacza orientację całej linii.
Ciągłość tego wyboru w rodzinach operatorów wynika
ze standardowej teorii wiązek wyznacznikowych:
orientację na końcu homotopii przenosimy
wzdłuż przedziału parametru do wyjściowego operatora.
Homotopia jest zgodna z działaniem cechowania.
Istotnie, na formach i funkcjach działanie
jest trywialne, a na spinorach mnożenie przez
$u$ jest zespolono-liniowe i przeplata
odpowiednie operatory Diraca. Na ich liniach
wyznacznikowych zachowuje orientację zespoloną.
Orientacja przenosi się więc na iloraz,
a nie tylko na przestrzeń konfiguracji.
Jeśli metryka się zmienia, wybieramy raz
maksymalną dodatnią podprzestrzeń $P$ formy
przecięcia i jej ortogonalne dopełnienie ujemne.
Rzut $H^2_+(X;\mathbb R)\to P$ wzdłuż
tego dopełnienia jest izomorfizmem:
jądro byłoby jednocześnie dodatnie i ujemne,
a obie przestrzenie mają wymiar $b_2^+$.
Pozwala to przenosić orientację homologiczną
spójnie między metrykami.
\end{proof}

\begin{stwierdzenie}[Przeliczanie punktów moduli]
\label{prop:tg-sw12-count}
Załóżmy $b_2^+(X)>0$ i $d(\mathfrak s)=0$.
Ustalmy orientację homologiczną oraz generyczne dane
$(g,\eta)$ bez rozwiązań redukowalnych. Wówczas
$\mathcal M_{\mathfrak s}(g,\eta)$ jest zbiorem skończonym,
a suma znaków jego punktów
$\sum_{[A,\psi]}\operatorname{sgn}([A,\psi])\in\mathbb Z$
jest dobrze określona. Przy $b_2^+>1$ nie zależy
od wyboru generycznych danych; przy $b_2^+=1$
należy dodatkowo ustalić komorę.
\end{stwierdzenie}

\begin{proof}
Z twierdzenia \ref{thm:tg-sw12-regular} przestrzeń moduli jest
zwartą gładką rozmaitością wymiaru zero, a więc
skończonym zbiorem punktów. Ich znaki wyznacza
orientacja linii wyznacznikowej. Przy zmianie danych
zachodzą dwie możliwości: gdy jednoparametrowa rodzina
unika redukowalnych, liczbę punktów zachowuje
jednowymiarowy zorientowany kobordyzm, dokładnie jak
w poprzednim rozdziale. Dla $b_2^+>1$ generyczna
ścieżka omija ścianę kowymiaru co najmniej dwa.
Dla $b_2^+=1$ ściana ma kowymiar jeden.
Łukowa spójność każdej komory i kobordyzm
wewnątrz niej zostały uzasadnione w rozdziale
\ref{ch:tg-sw-analiza}, po stwierdzeniu
\ref{prop:tg-sw-wall}. Liczba może się zmienić
przy przekroczeniu ściany.
\end{proof}

\section{Lokalny model ściany i rachunek zmiany komory}
\label{sec:tg-sw12-wall}

\begin{stwierdzenie}[Indeks w punkcie redukowalnym]
\label{prop:tg-sw12-wall-index}
Niech $b_1(X)=0$, $b_2^+(X)=1$, a
$d(\mathfrak s)=0$. Wtedy
$\operatorname{ind}_\mathbb C D_A=1$
dla każdego połączenia $A$ zgodnego z $\mathfrak s$.
\end{stwierdzenie}

\begin{proof}
Wzór na indeks operatora deformacyjnego ma w tym
przypadku postać
$0=d(\mathfrak s)=2\operatorname{ind}_\mathbb C D_A
+b_1-1-b_2^+=2\operatorname{ind}_\mathbb C D_A-2$.
Stąd żądana równość. Sam indeks nie mówi jeszcze,
że jądro Diraca ma wymiar zespolony jeden: wymagałoby
to także zanikania jego cokernelu.
\end{proof}

\begin{twierdzenie}[Lokalny mechanizm zmiany komory]
\label{thm:tg-sw12-wall-model}
Przy założeniach stwierdzenia
\ref{prop:tg-sw12-wall-index} wybierzmy
jednoparametrową rodzinę danych, regularną
na części nieredukowalnej, z regularnymi
danymi końcowymi, przecinającą ścianę
redukowalnych poprzecznie w jednym punkcie.
Załóżmy, że w tym punkcie operator Diraca jest
surjektywny. Wówczas lokalna przestrzeń
moduli parametrycznych po ilorazowaniu
przez stabilizator $S^1$ jest homeomorficzna
z otoczeniem zera w półprostej.
W konsekwencji
$|SW_X^+(\mathfrak s)-SW_X^-(\mathfrak s)|=1$.
Nie wymagamy niezerowości współczynnika
kwadratowego w parametrze spinorowym.
\end{twierdzenie}

\begin{proof}
Przy $b_1=0$ na ścianie jest dokładnie jedna
klasa redukowalna (opis po stwierdzeniu
\ref{prop:tg-sw-wall}). Oznaczmy jej reprezentanta
przez $(A,0)$ i parametr przejścia przez $\tau=0$.
Używamy lokalnego przekroju połączeń
$\mathrm d_g^*a=0$ względem metryki w tym punkcie,
usuwając transformacje cechowania zadane
przez funkcje o średniej zero.
Przy $b_1=0$ każda transformacja cechowania
ma postać $u=e^{if}$ dla rzeczywistej funkcji $f$
(podniesienie do nakrycia $\mathbb R\to S^1$,
gdy klasa całkowita z
\ref{prop:tg-sw12-lattice} jest zerowa).
Dla dwóch połączeń w przekroju Coulomba
związanych taką transformacją mamy
$\mathrm d^*\mathrm df=0$; na spójnej,
zamkniętej rozmaitości $f$ jest stała.
Z kolei dowolne pobliskie połączenie można
sprowadzić do tego przekroju, rozwiązując
równanie Poissona z warunkiem zerowej średniej,
jak w dowodzie twierdzenia
\ref{thm:tg-sw-compact}.
Pozostaje zatem dokładnie działanie
stałych faz $S^1$.
W lokalnych identyfikacjach wiązek dla metryk
rodziny linearyzacja w tym przekroju ma w
punkcie redukowalnym postać
\[
(a,\phi,\dot\tau)\longmapsto
(D_A\phi,\ \mathrm d_g^+a+\dot\tau\,\beta),
\]
gdzie $\beta$ jest wariacją równania krzywizny
w kierunku parametru, z wartościami w
$i\Lambda^2_+$. Człon mieszany w równaniu
Diraca znika, gdyż spinor wyjściowy jest zerowy;
z tego samego powodu znika różniczka $q$.

Poprzeczność do ściany oznacza, że rzut
$\beta/i$ na $\mathcal H_g^+$ jest niezerowy.
Istotnie, różniczkując harmoniczne parowanie
równania krzywizny, człon pochodzący
z różniczkowania samego parowania znika
w punkcie, w którym równanie jest spełnione.
Operator $\mathrm d_g^+$ na przekroju Coulomba
ma zerowe jądro, gdyż $b_1=0$,
i cokernel $i\mathcal H_g^+$ wymiaru jeden.
Dodanie kierunku parametru daje zatem
izomorfizm tej części na przestrzeń
samodualnych form w odpowiednich normach Sobolewa.
Operator Diraca jest surjektywny z założenia,
a z indeksu wynika $\ker D_A\cong\mathbb C$.

Wybieramy dopełnienie $\ker D_A$ w przestrzeni
spinorów, niezmiennicze względem stałych faz,
np. dopełnienie ortogonalne w $L^2$.
Różniczka względem zmiennych połączeniowych,
parametru i tego dopełnienia jest izomorfizmem.
Twierdzenie o funkcji uwikłanej daje lokalny
gładki model rozwiązań parametrycznych
przed ilorazowaniem, ze współrzędną
$z\in\ker D_A\cong\mathbb C$.
Jednoznaczność rozwiązań w dopełnieniu
zapewnia zgodność modelu z działaniem faz
$z\mapsto uz$, $u\in S^1$.
Dla $z=0$ otrzymujemy konfigurację redukowalną,
a dla $z\ne0$ spinor ma niezerowy rzut
na $\ker D_A$ i jest nieredukowalny.
Model lokalnego ilorazu przez stałe fazy
to zatem $\mathbb C/S^1\cong[0,\infty)$,
z parametrem $r=|z|^2$.

Po usunięciu małego otoczenia punktu
redukowalnego pozostaje zwarty zorientowany
kobordyzm jednowymiarowy. Ma on dodatkowy
punkt brzegu $r=\delta>0$, niezależnie od
tego, czy projekcja lokalnej półprostej na
parametr $\tau$ jest monotoniczna.
Wkład tego punktu do algebraicznej liczby
punktów brzegu jest równy $1$ lub $-1$.
Stąd otrzymujemy wzór na wartość
bezwzględną skoku niezmiennika.
W argumentacji wykorzystujemy standardowy
lokalny przekrój przy redukowalnych i teorię
parametrycznych orientacji, jak w twierdzeniu
\ref{thm:tg-sw-wallcross}. Generyczna
surjektywność Diraca przy przejściu przez
ścianę jest wynikiem analitycznym przyjętym
w tamtym twierdzeniu, a nie wnioskiem
z samej równości indeksowej.
\end{proof}

\begin{center}
\begin{tikzpicture}[>=Stealth,scale=.9]
  \draw[tgNavy!60,->] (-4.2,0)--(4.25,0)
       node[right,font=\scriptsize] {$\tau$};
  \draw[tgNavy!60,->] (0,-.15)--(0,2.5)
       node[above,font=\scriptsize] {$r=|z|^2$};
  \draw[tgTeal,line width=2.6pt] (0,0) .. controls
       (-.15,.8) and (-1.5,1.6) .. (-2.9,2.2);
  \fill[tgCoral] (0,0) circle (.09);
  \node[text=tgCoral,font=\small,anchor=west] at (.23,.36)
       {punkt redukowalny};
  \node[text=tgTeal,font=\small] at (-2.55,2.55)
       {$\tau=h(r)$};
  \draw[tgPurple,dashed] (-3.4,1.25)--(2.7,1.25);
  \node[text=tgPurple,font=\scriptsize,anchor=west] at (.35,1.43)
       {przekrój $r=\delta$};
\end{tikzpicture}

{\footnotesize Przykładowy lokalny model ilorazu przez $S^1$.
Istotny jest jeden punkt brzegu po odcięciu otoczenia
rozwiązania redukowalnego. Rysunek nie zakłada
niezerowości pochodnej funkcji $h$ w zerze.}
\end{center}

\section{Co zostało dowiedzione, a co wymaga dalszej teorii?}
\label{sec:tg-sw12-scope}

W tej części dowiedliśmy jawnie tożsamości algebraicznych
dla $q$, dokładności symbolu kompleksu samodualnego,
kontroli pełnej krzywizny i wzorów na indeks
\emph{przy założeniu} twierdzenia Atiyaha--Singera.
Pokazaliśmy także krok po kroku, jak oszacowanie spinora,
Hodge'owskie ustalenie cechowania i bootstrap eliptyczny
prowadzą do zwartości. Dla ustalonej metryki generyczną regularność
wyprowadziliśmy z jednoznacznego przedłużenia,
zwartości i skończeniewymiarowego twierdzenia Sarda.

Nie udowodniliśmy w tym skrypcie od podstaw:
(i) pełnej teorii Fredholma w skali Sobolewa,
(ii) globalnych oszacowań eliptycznych i twierdzeń o
zwartości zanurzeń Sobolewa,
(iii) twierdzenia indeksowego Atiyaha--Singera,
(iv) parametrycznej wersji Sarda--Smale'a,
(v) generycznej surjektywności operatora Diraca
w punkcie redukowalnym oraz pełnej teorii orientacji
linii wyznacznikowych w rodzinach.
Sam lokalny model przy założonej surjektywności
Diraca został wyprowadzony z twierdzenia o funkcji
uwikłanej; nie przyjmujemy już dodatkowego
twierdzenia o niezerowości współczynnika
kwadratowego. Dowody ogólnych twierdzeń analitycznych
należą do analizy globalnej i teorii cechowania.

\begin{uwaga}[Most do następnej iteracji]
W następnych rozdziałach zastosujemy ten aparat
do konstrukcji Parka: oprócz formy przecięcia
sprawdzimy $\pi_1$, typ topologiczny oraz niezmienniki
Seiberga--Wittena w odpowiedniej komorze.
\end{uwaga}


% ========================================================================
% Iteracja 13: konstrukcja Parka. Umiescic po analizie SW (iteracja 12),
% przed rozdzialem o nowoczesnych niezmiennikach wezlow.
% ========================================================================
